\chapter{Минимальный родитель обменного моста и знак действия}

Предыдущая проверка оставила две возможности: включить гладкий компактный идеал
Тёплица в координатную алгебру либо считать мост нечётным эндоморфизмом
градуированного соответствия. Настоящая проверка выбирает минимальный
кинематический родитель и сразу выясняет, поддерживает ли его каноническое
положительное действие рождение дефекта.

\section{Сверка с проектом и литературой}

Том V уже запретил объявлять полный операторный контейнер координатной
алгеброй: это расширяет внутренние флуктуации и калибровочные направления.
Одновременно он допустил семейный нечётный эндоморфизм как самостоятельные
данные градуированного соответствия, не изменяющие алгебру Стандартной
модели.

Это разделение соответствует литературе. Суперсвязность на
$\mathbb Z_2$-градуированном модуле может содержать нечётную нуль-форменную
эндоморфную часть; см. \path{arXiv:math-ph/9911006} и
\path{arXiv:0810.0820}. В неограниченной бивариантной $K$-теории оператор,
модуль и связность являются данными соответствия, а не обязаны возникать
из расширения координатной алгебры; см. \path{arXiv:0904.4383} и
\path{arXiv:1911.05008}.

\section{Почему компактная координатная ветвь не минимальна}

Пусть гладкая координатная алгебра содержит все конечные матричные единицы
$E_{ij}$ на пространстве Харди. Для оператора числа $N$ выполнено
\begin{equation}
 [N,E_{kj}]=(k-j)E_{kj}.
\end{equation}
Если $k\ne j$, то
\begin{equation}
 \frac1{k-j}E_{ik}[N,E_{kj}]=E_{ij}.
 \label{eq:v6-all-compacts-as-oneforms}
\end{equation}
Следовательно, представленное пространство одноформ содержит не только
$P_0=E_{00}$, но все конечные матричные единицы. При увеличении обрезания
оно порождает весь гладкий компактный идеал.

\begin{proposition}[Неизолированность компактного моста]
Добавление конечноматричного идеала как координатной алгебры не вводит одну
амплитуду обменного моста. Оно открывает полный компактный сектор одноформ.
Без дополнительного селектора или меры этот вариант не является
минимальным родителем искомого механизма.
\end{proposition}

Это не математический запрет расширения Тёплица. Это запрет считать его
экономным физическим выводом одного поля.

\section{Real-суперсвязностный родитель}

Рассмотрим две сопряжённые ориентации как вещественный модуль над
$M_{15}(\mathbb C)_{\mathbb R}$. Между стандартным и сопряжённым
представлениями существует двумерное над $\mathbb R$ пространство
интертвинеров, соответствующее одному комплексному скаляру. Условие
обменной вещественности
\begin{equation}
 RT_\lambda=T_\lambda^*R
\end{equation}
удаляет мнимую фазу и оставляет одну вещественную амплитуду $\lambda$.
Каноническая нуль-форменная часть суперсвязности имеет вид
\begin{equation}
 \Phi_\lambda=\lambda
 \begin{pmatrix}0&P_q\\P_q&0\end{pmatrix}.
 \label{eq:v6-real-bridge-superconnection-field}
\end{equation}

Она действует на ориентационном тёплицевом соответствии, а не на трёх
заряженных рёбрах $u,d,e$. Поэтому координатная алгебра Стандартной модели
и ранее замороженный модуль $\mathcal Y_\rho$ не расширяются.

\begin{theorem}[Минимальный кинематический родитель]
Среди двух проверенных вариантов минимальным является обменный мост как
вещественный нечётный эндоморфизм градуированного соответствия. Он вводит
одну вещественную амплитуду, сохраняет координатную алгебру и не добавляет
четвёртого физического ребра в $H_{15}$.
\end{theorem}

Теорема является кинематической. Она ещё не утверждает, что действие
выбирает ненулевой материальный дефект.

\section{Знак канонического положительного действия}

Для семейства $T_\lambda$ предыдущей проверки точные тождества дают
\begin{align}
 I-T_\lambda^*T_\lambda&=(1-\lambda^2)P_q
 \quad\text{на одной ветви},\\
 I-T_\lambda T_\lambda^*&=(1-\lambda^2)P_q
 \quad\text{на сопряжённой ветви}.
\end{align}
Самое экономное положительное действие кривизны поэтому равно
\begin{equation}
 \mathcal S_{\rm can}(\lambda)
 =\tau\!\left((I-T_\lambda^*T_\lambda)^2
 +(I-T_\lambda T_\lambda^*)^2\right)
 =\frac{2}{7}(1-\lambda^2)^2.
 \label{eq:v6-canonical-bridge-action}
\end{equation}
Отсюда
\begin{align}
 \mathcal S_{\rm can}(0)&=\frac27,&
 \mathcal S_{\rm can}(1)&=0,\\
 \mathcal S_{\rm can}''(0)&=-\frac87,&
 \mathcal S_{\rm can}''(1)&=\frac{16}{7}.
\end{align}

Таким образом, дефектное состояние $\lambda=0$ является неустойчивой
вершиной, а полностью замкнутый мост $\lambda=1$ --- устойчивым минимумом.
Естественная положительная норма стремится уничтожить пару нулевых мод, а
не родить её.

\begin{nogocriterion}[Положительная кривизна не рождает материю]
Нельзя объявлять суперсвязностный мост механизмом рождения материи только
потому, что он даёт единое пространство конфигураций. Каноническое
положительное действие выбирает обратимый замкнутый представитель. Для
рождения дефекта необходимо независимо вывести член или граничное условие,
дестабилизирующее $\lambda=1$; менять знак квадрата вручную запрещено.
\end{nogocriterion}

\section{Статус механизма}

Обменная суперсвязность закрывает вопрос минимального математического
носителя, но даёт отрицательный ответ на первый динамический тест. Теперь
искать следует не ещё одну кинематику, а уже существующий в проекте источник
дестабилизации замкнутого моста: фермионный детерминант, относительную
эта-фазу, пространственное граничное условие или топологически вынужденный
ноль. Каждый кандидат обязан изменить гессиан при $\lambda=1$ без нового
ручного коэффициента.

\begin{protocol}
\begin{enumerate}
 \item компактный проектор как представленная одноформа: \textbf{да};
 \item изоляция одного проектора после расширения алгебры:
 \textbf{нет, возникает весь компактный сектор};
 \item компактное расширение как минимальный физический родитель:
 \textbf{не проходит};
 \item Real-линейный нечётный эндоморфизм соответствия:
 \textbf{проходит кинематически};
 \item число вещественных амплитуд после обмена: \textbf{одна};
 \item изменение $H_{15}$ и рёбер $u,d,e$: \textbf{нет};
 \item каноническое положительное действие: \textbf{выведено};
 \item устойчивость дефектной точки $\lambda=0$: \textbf{нет};
 \item устойчивость замкнутого моста $\lambda=1$: \textbf{да};
 \item спонтанное рождение материи этим действием: \textbf{не получено};
 \item следующая проверка:
 \textbf{поиск дестабилизации замкнутого моста в доказанном багаже проекта}.
\end{enumerate}
\end{protocol}

Машинный сертификат сохранён в
\path{s2t_v6_exchange_bridge_minimal_parent_gate_results.json}.